\documentclass[runningheads]{llncs}

\usepackage[T1]{fontenc}
\usepackage{graphicx}
\usepackage{amsmath}
\usepackage{amssymb}
\usepackage{caption}
\usepackage{subcaption}
\usepackage{booktabs}
\usepackage{multirow}
\usepackage{hyperref}

\graphicspath{{figures/}}

\begin{document}

\title{Pattern Layout Optimization for Seam-Aware Texture Alignment \\
       Supplementary Material}

\author{Anonymous}
\authorrunning{Anonymous}
\institute{Anonymous}

\maketitle

% ---------------------------------------------------------------------------
\section{Full Parameter Specification}
\label{sup:params}

Table~\ref{tab:hyperparams} lists every hyperparameter used in the main
experiments. All reported results use the same configuration except
where explicitly noted (e.g.\ the $\text{GA}_{\text{no-}\kappa}$
ablation sets $p_\kappa = 0$; the B3 random-search baseline uses the
same bounds but no operators).

\begin{table}[h]
\centering
\caption{Full hyperparameter specification.}
\label{tab:hyperparams}
\setlength{\tabcolsep}{5pt}
\begin{tabular}{llc}
\toprule
Group & Parameter & Value \\
\midrule
\multirow{7}{*}{Search space}
 & $\boldsymbol{\delta}$ bounds  $[\delta_\text{lo}, \delta_\text{hi}]$ & $[0.2, 0.8]$ \\
 & $\boldsymbol{\delta}$ dimension (onesie) & $2m = 2 \times 13$ \\
 & $\boldsymbol{\delta}$ dimension (pants)  & $2m = 2 \times 5$ \\
 & Phase bins $K$                           & $8$ \\
 & Grain orientations $|\{\rho\}|$          & $4$ (fixed to $0$ in experiments) \\
 & Nesting heuristics $|\mathcal{H}|$       & $3$ (fixed to bottom-left) \\
 & Permutation $\pi$                        & area-sorted (fixed) \\
\midrule
\multirow{8}{*}{GA operators}
 & Population size                          & $100$ \\
 & Generations                              & $20$ \\
 & Elite count                              & $4$ \\
 & Tournament size $k$                      & $4$ \\
 & Crossover probability $p_c$              & $0.7$ \\
 & Mutation probability $p_m$               & $0.7$ \\
 & $\boldsymbol{\delta}$ Gaussian $\sigma_\delta$ & $0.20$ \\
 & $\boldsymbol{\kappa}$ per-gene flip $p_\kappa$ & $0.35$ \\
\midrule
\multirow{4}{*}{Fitness weights}
 & $w_\text{fabric}$                        & $1.0$ \\
 & $w_\text{align}$                         & $2.5$ \\
 & $w_\text{fit}$                           & $10.0$ \\
 & Tiebreaker (selection)                   & Pareto dominance, sum fallback \\
\midrule
\multirow{6}{*}{Environment}
 & Fabric roll width $W$                    & $1{,}500$ mm \\
 & Texture period $(p_u, p_v)$              & $(50, 50)$ mm \\
 & Body counts $N$                          & $\{5, 10, 25, 50, 100\}$ \\
 & Seeds per condition                      & $10$ \\
 & Mesh                                     & SMPL female, neutral pose \\
 & Wallpaper group (main results)           & PM (horizontal stripes) \\
\bottomrule
\end{tabular}
\end{table}

% ---------------------------------------------------------------------------
\section{Runtime Breakdown}
\label{sup:runtime}

Table~\ref{tab:runtime} decomposes the per-evaluation cost of a single
individual on a single CPU core. Geometry dominates: LSCM flattening
and geodesic cutting account for $\sim 90\%$ of the time, Stage~2
Levenberg--Marquardt is a small constant, and the nesting engine itself
is negligible relative to the rest.

\begin{table}[h]
\centering
\caption{Per-evaluation runtime breakdown, single CPU core.}
\label{tab:runtime}
\setlength{\tabcolsep}{5pt}
\begin{tabular}{lc}
\toprule
Component & Time per evaluation \\
\midrule
Geometry pipeline (cut + LSCM flatten)     & ${\sim}$4.5\,s \\
Stage 2 Levenberg--Marquardt                & ${\sim}$0.1\,s \\
Nesting (bottom-left, STRtree collision)    & $<$0.01\,s \\
\midrule
Total per individual                        & ${\sim}$5\,s \\
\midrule
Full GA run (pop $=100$, gens $=20$, $2{,}000$ evals)     & ${\sim}$3 h \\
Full B3 run ($2{,}000$ random samples)                     & ${\sim}$3 h \\
B2 onesie ($5 \times 10^5$ $\boldsymbol{\kappa}$ samples, geometry run once) & ${\sim}$3 h \\
\bottomrule
\end{tabular}
\end{table}

% ---------------------------------------------------------------------------
\section{Geometry Failure Rates}
\label{sup:failures}

A small fraction of candidate $\boldsymbol{\delta}$ configurations
produce degenerate mesh cuts or LSCM flattening failures. These
individuals receive a penalty fitness and are discarded by selection,
so they do not directly affect the reported $F_\Sigma$ values but do
reduce the effective population size. Table~\ref{tab:failures} reports
the failure rate per body count.

\begin{table}[h]
\centering
\caption{Fraction of evaluations that fail in the geometry pipeline
(geodesic or LSCM failure) per body count, aggregated over all seeds.
$^\dagger$Placeholder; to be filled from run logs.}
\label{tab:failures}
\begin{tabular}{cc}
\toprule
$N$ & Failure rate$^\dagger$ \\
\midrule
5    & -- \\
10   & -- \\
25   & -- \\
50   & -- \\
100  & -- \\
\bottomrule
\end{tabular}
\end{table}

% ---------------------------------------------------------------------------
\section{Per-Seed Convergence Trajectories}
\label{sup:convergence}

Figure~\ref{fig:conv_per_seed} overlays all 10 seeds per body count,
showing the raw step-like structure that the main-paper
Figure~\ref{fig:convergence} summarises as a median and interquartile
range. The individual trajectories clearly show that the step positions
vary seed by seed, consistent with the Voronoi interpretation:
different random initialisations land in different plateaus, and
improvements occur whenever a mutation happens to cross into a
better-flattening region of the landscape.

% TODO: add per-seed convergence figure once 10-seed data is ready.
% \begin{figure}[h]
% \centering
% \includegraphics[width=\linewidth]{convergence_per_seed}
% \caption{Per-seed convergence trajectories, all 10 seeds overlaid per
%   body count. Horizontal axis: generation. Vertical axis: best-so-far
%   $F_\Sigma$.}
% \label{fig:conv_per_seed}
% \end{figure}

% ---------------------------------------------------------------------------
\section{Fit Guard $f_\text{fit}$ Distribution}
\label{sup:ffit}

Figure~\ref{fig:ffit_dist} shows the empirical distribution of
$f_\text{fit}$ across all 10 seeds and 5 body counts. The guard is
zero in the majority of runs, confirming that the $w_\text{fit}=10$
penalty is effective at preventing garment shrinkage during
$\boldsymbol{\delta}$ search.

% TODO: add f_fit histogram figure once 10-seed data is ready.
% \begin{figure}[h]
% \centering
% \includegraphics[width=0.6\linewidth]{ffit_distribution}
% \caption{Histogram of $f_\text{fit}$ values across all seeds and body
%   counts. The mass at zero confirms the guard constraint effectively
%   prevents shrinkage.}
% \label{fig:ffit_dist}
% \end{figure}

% ---------------------------------------------------------------------------
\section{Full Pants Results (All $N$)}
\label{sup:pants_full}

% TODO: fill full pants table for all N in {5, 10, 25, 50, 100}.

% ---------------------------------------------------------------------------
\section{Sensitivity to Phase Resolution $K$}
\label{sup:k_sensitivity}

To assess how the GA's performance scales with the phase bin
resolution, we run the full GA on pants (where B2 is exact) with
$K \in \{4, 8, 16\}$ and 10 seeds per configuration. Finer phase bins
give Stage 2 and the GA finer-grained control over the residual
phase mismatch.

% TODO: add K sensitivity table.

% ---------------------------------------------------------------------------
\section{Population Diversity}
\label{sup:population}

Figure~\ref{fig:population} visualises an entire population in a
single generation, showing the diversity of seam placements and
resulting nesting layouts across individuals. This demonstrates
qualitatively that the GA maintains a broad search distribution
throughout the run rather than collapsing to a single mode.

% TODO: add population figure from results/experiments/onesie_sleeves_v2/.

% ---------------------------------------------------------------------------
\section{Additional Wallpaper Group Examples}
\label{sup:wallpaper_examples}

For each of the eight supported wallpaper groups, Figure~\ref{fig:wg_examples}
shows a high-quality (best individual) and a low-quality (worst
individual) example from a single GA run, illustrating the visual
range of alignment outcomes the optimiser navigates during search.

% TODO: add good/bad wallpaper examples once wallpaper runs complete.

% ---------------------------------------------------------------------------
\section{Results on Additional Body Shapes}
\label{sup:body_shapes}

To verify that the GA's performance generalises across different body
shapes, we repeat the main onesie experiment on two additional SMPL
female meshes with substantially different $\beta$ shape parameters.

% TODO: add additional body shape results if compute allows.

% ---------------------------------------------------------------------------
\section{Results on Shirts}
\label{sup:shirts}

The sleeveless shirt (\texttt{upper}, $M=2$) is a simpler garment
type not included in the main paper; we report it here for
completeness.

% TODO: fill sleeveless shirt results.

\end{document}
